\documentclass{article}
\usepackage{amsmath}
\usepackage{amssymb}

\title{Tilde Engine Benchmark & Stress Test}
\author{Google Antigravity & Moth}
\date{\today}

\begin{document}

\maketitle

\section{Introduction to Tilde}
This document serves as a live rendering benchmark for \textbf{Tilde}, a native C++17 dual-pane Markdown and \LaTeX\ editor built on top of Qt 6 and KDE Frameworks (\texttt{KF6::TextEditor}). Unlike modern Electron or web-based text editors that consume over 300~MB of memory just to render basic rich text, Tilde leverages native system APIs (\texttt{QPdfView} and \texttt{QProcess}) to maintain an ultra-lightweight memory footprint.

\section{Mathematical Expressions}
Below are several mathematical equations rendered live in the background by \texttt{pdflatex} and displayed via \texttt{Qt6::PdfWidgets}:

\subsection{The Quadratic Formula}
The roots of a second-degree polynomial $ax^2 + bx + c = 0$ are given by:
\[
x = \frac{-b \pm \sqrt{b^2 - 4ac}}{2a}
\]

\subsection{Gaussian Integral}
A cornerstone of probability theory and quantum mechanics:
\[
\int_{-\infty}^{\infty} e^{-x^2} \, dx = \sqrt{\pi}
\]

\subsection{Maxwell's Equations in Differential Form}
In a vacuum where no charges or currents are present, Maxwell's equations take the following symmetric form:
\begin{align}
\nabla \cdot \mathbf{E} &= 0 \\
\nabla \cdot \mathbf{B} &= 0 \\
\nabla \times \mathbf{E} &= -\frac{\partial \mathbf{B}}{\partial t} \\
\nabla \times \mathbf{B} &= \mu_0 \epsilon_0 \frac{\partial \mathbf{E}}{\partial t}
\end{align}

\section{Conclusion}
If you can see this document rendered cleanly with crisp vector typography and mathematical symbols on the right half of your screen, the asynchronous \texttt{QProcess} and debounced \texttt{QTimer} compilation pipeline is functioning with 100\% precision.

\end{document}
